% ============================================================
\chapter{Text Generation and Decoding}
\label{ch:generation}
% ============================================================

A language model that predicts the next word is, quite literally, a text generator in waiting. But between the distribution $P(x_t \mid x_{<t})$ and coherent, fluent, informative text lies a chasm --- that of the \emph{decoding strategy}. Systematically choosing the most probable word (greedy decoding) produces repetitive, dull text. Na\"ively sampling from the distribution produces incoherent text. The solutions --- beam search, top-$k$, nucleus sampling (top-$p$), temperature --- are all trade-offs between \emph{quality} and \emph{diversity}. Ari Holtzman and collaborators, in their paper ``The Curious Case of Neural Text Degeneration'' (2020), showed that nucleus sampling elegantly solves the degeneration problem. This chapter explores these strategies and their probabilistic foundations.

\begin{intuition}[title=\textbf{From language model to generator}]
An autoregressive language model defines $P(x_t \mid x_{<t})$. Generating text
amounts to sequentially sampling from this distribution. The choice of decoding
strategy radically influences the quality, diversity, and coherence of the
produced text.
\end{intuition}

% ------------------------------------------------------------
\section{Greedy Decoding}
\label{sec:gen:greedy}
% ------------------------------------------------------------

\begin{definition}[Greedy Decoding]
Greedy decoding selects the most probable token at each step:
\[
x_t = \argmax_{w \in \mathcal{V}} P(w \mid x_{<t})
\]
\end{definition}

\begin{warning}[title=\textbf{Limitations of greedy decoding}]
Greedy decoding does not guarantee the globally optimal sequence. It can produce
repetitions and lacks diversity. The most probable sequence is not always the
most natural one.
\end{warning}

% ------------------------------------------------------------
\section{Beam Search}
\label{sec:gen:beam}
% ------------------------------------------------------------

\begin{definition}[Beam Search]
Beam search maintains the $B$ most probable partial sequences (\emph{beams})
at each decoding step:
\[
\mathcal{B}_t = \text{top-}B\left\{(s \oplus w, \text{score}(s) + \log P(w \mid s)) : s \in \mathcal{B}_{t-1}, w \in \mathcal{V}\right\}
\]
where $B$ is the beam width.
\end{definition}

\begin{proposition}[Properties of Beam Search]
\begin{enumerate}
    \item For $B = 1$: equivalent to greedy decoding.
    \item For $B = \abs{\mathcal{V}}^T$: equivalent to exhaustive search.
    \item In practice, $B \in [4, 10]$ offers a good trade-off.
\end{enumerate}
\end{proposition}

\subsection{Length Normalization}

Beam search favors shorter sequences (more $\log P < 0$ terms for longer ones).
We normalize the score by length:
\[
\text{score}_{\text{norm}}(s) = \frac{1}{\abs{s}^\alpha} \sum_{t=1}^{\abs{s}} \log P(x_t \mid x_{<t})
\]
where $\alpha \in [0.6, 1.0]$ is a hyperparameter (Wu et al., 2016).

% ------------------------------------------------------------
\section{Stochastic Sampling}
\label{sec:gen:sampling}
% ------------------------------------------------------------

\begin{definition}[Temperature Sampling]
Temperature sampling with $\tau > 0$ modifies the distribution:
\[
P_\tau(w \mid x_{<t}) = \frac{\exp(\text{logit}_w / \tau)}{\sum_{w'} \exp(\text{logit}_{w'} / \tau)}
\]
\begin{itemize}
    \item $\tau \to 0$: converges to greedy decoding.
    \item $\tau = 1$: original distribution.
    \item $\tau > 1$: more uniform distribution (more diversity).
\end{itemize}
\end{definition}

\begin{definition}[Top-$k$ Sampling]
Top-$k$ sampling (Fan et al., 2018) restricts sampling to the $k$ most probable
tokens:
\[
P_{\text{top-}k}(w \mid x_{<t}) = \begin{cases}
\frac{P(w \mid x_{<t})}{\sum_{w' \in \mathcal{V}_k} P(w' \mid x_{<t})} & \text{if } w \in \mathcal{V}_k \\
0 & \text{otherwise}
\end{cases}
\]
where $\mathcal{V}_k$ contains the $k$ most probable tokens.
\end{definition}

\begin{definition}[Top-$p$ Sampling (Nucleus)]
Top-$p$ sampling (Holtzman et al., 2020) selects the smallest set
$\mathcal{V}_p$ such that:
\[
\sum_{w \in \mathcal{V}_p} P(w \mid x_{<t}) \geq p
\]
then samples from $\mathcal{V}_p$ (after renormalization).
\end{definition}

\begin{bestpractice}[title=\textbf{Combining strategies}]
In practice, temperature, top-$k$, and top-$p$ are often combined. For example,
$\tau = 0.7$, $k = 50$, $p = 0.9$ offers a good balance between coherence and
diversity for creative generation.
\end{bestpractice}

% ------------------------------------------------------------
\section{Repetition Penalties}
\label{sec:gen:repetition}
% ------------------------------------------------------------

\begin{definition}[Repetition Penalty]
To avoid repetitions, already-generated tokens are penalized:
\[
\text{logit}'_w = \begin{cases}
\text{logit}_w / \theta & \text{if } w \in x_{<t} \text{ and } \text{logit}_w > 0 \\
\text{logit}_w \cdot \theta & \text{if } w \in x_{<t} \text{ and } \text{logit}_w \leq 0
\end{cases}
\]
where $\theta > 1$ is the penalty factor.
\end{definition}

% ------------------------------------------------------------
\section{Contrastive Decoding}
\label{sec:gen:contrastive}
% ------------------------------------------------------------

\begin{definition}[Contrastive Decoding]
Contrastive decoding (Li et al., 2023) selects tokens that maximize the
log-probability difference between an expert model $\pi_{\text{exp}}$ and an
amateur model $\pi_{\text{ama}}$:
\[
x_t = \argmax_{w \in \mathcal{V}_p} \left[\log \pi_{\text{exp}}(w \mid x_{<t}) - \alpha \log \pi_{\text{ama}}(w \mid x_{<t})\right]
\]
\end{definition}

% ------------------------------------------------------------
\section{Speculative Decoding}
\label{sec:gen:speculative}
% ------------------------------------------------------------

\begin{definition}[Speculative Decoding]
Speculative decoding (Leviathan et al., 2023) uses a small fast model
(\emph{draft model}) to propose $K$ tokens, then a large model verifies them in
parallel. Accepted tokens are those whose probability under the large model is
sufficient, accelerating generation without changing the output distribution.
\end{definition}

% ------------------------------------------------------------
\section{Generation Quality Metrics}
\label{sec:gen:metrics}
% ------------------------------------------------------------

\begin{keyformulas}[title=\textbf{Generation metrics}]
\begin{align}
\text{BLEU} &= \text{BP} \cdot \exp\!\left(\sum_{n=1}^{N} w_n \log p_n\right) & \text{($n$-gram precision)} \\
\text{ROUGE-L} &= F_1(\text{LCS}) & \text{(longest common subsequence)} \\
\text{METEOR} &= F_1 \cdot (1 - \text{Penalty}) & \text{(flexible alignment)} \\
\text{BERTScore} &= F_1(\text{BERT cosine sim}) & \text{(contextual similarity)}
\end{align}
\end{keyformulas}

% ------------------------------------------------------------
\section{Implementation}
\label{sec:gen:implementation}
% ------------------------------------------------------------

\begin{codeblock}[title=\textbf{Decoding strategies with Hugging Face}]
\begin{minted}{python}
from transformers import AutoTokenizer, AutoModelForCausalLM

model_name = "gpt2-medium"
tokenizer = AutoTokenizer.from_pretrained(model_name)
model = AutoModelForCausalLM.from_pretrained(model_name)

prompt = "The future of artificial intelligence"
input_ids = tokenizer.encode(prompt, return_tensors="pt")

# Greedy decoding
greedy = model.generate(input_ids, max_length=50)

# Beam search
beam = model.generate(input_ids, max_length=50, num_beams=5,
                      length_penalty=0.8, no_repeat_ngram_size=3)

# Top-k + Top-p sampling
sampled = model.generate(input_ids, max_length=50, do_sample=True,
                         temperature=0.7, top_k=50, top_p=0.9,
                         repetition_penalty=1.2)

for name, ids in [("Greedy", greedy), ("Beam", beam), ("Sampled", sampled)]:
    print(f"{name}: {tokenizer.decode(ids[0], skip_special_tokens=True)}")
\end{minted}
\end{codeblock}

\begin{codeblock}[title=\textbf{Step-by-step decoding}]
\begin{minted}{python}
import torch
import torch.nn.functional as F

def generate_top_p(model, input_ids, max_new_tokens=50, temperature=0.8,
                   top_p=0.9):
    generated = input_ids.clone()
    for _ in range(max_new_tokens):
        with torch.no_grad():
            outputs = model(generated)
            logits = outputs.logits[:, -1, :] / temperature

        # Top-p filtering
        sorted_logits, sorted_indices = torch.sort(logits, descending=True)
        cumulative_probs = torch.cumsum(
            F.softmax(sorted_logits, dim=-1), dim=-1
        )
        mask = cumulative_probs - F.softmax(sorted_logits, dim=-1) >= top_p
        sorted_logits[mask] = float('-inf')
        logits.scatter_(1, sorted_indices, sorted_logits)

        probs = F.softmax(logits, dim=-1)
        next_token = torch.multinomial(probs, num_samples=1)
        generated = torch.cat([generated, next_token], dim=-1)

        if next_token.item() == tokenizer.eos_token_id:
            break
    return generated
\end{minted}
\end{codeblock}

% ------------------------------------------------------------
\section{Controlled Generation}
\label{sec:gen:controlled}
% ------------------------------------------------------------

\begin{definition}[Controlled Generation]
\emph{Controlled generation} aims to steer a language model's output according
to specified attributes (tone, topic, style) without retraining. Approaches:
\begin{itemize}
    \item \textbf{PPLM} (Dathathri et al., 2020): gradient perturbation.
    \item \textbf{Classifier-free guidance}: interpolation between conditional
          and unconditional distributions.
    \item \textbf{Constrained decoding}: hard lexical constraints.
\end{itemize}
\end{definition}

% ------------------------------------------------------------
\section{Exercises}
\label{sec:gen:exercises}
% ------------------------------------------------------------

\begin{exercise}[$\star$]
Show that greedy decoding is a special case of beam search with $B = 1$.
\end{exercise}

\begin{exercise}[$\star$]
Compute the top-$p$ distribution for $p = 0.9$ given the probabilities
$(0.4, 0.3, 0.15, 0.1, 0.05)$.
\end{exercise}

\begin{exercise}[$\star\star$]
Implement beam search with length normalization. Test with different values of
$\alpha$ and $B$ on GPT-2.
\end{exercise}

\begin{exercise}[$\star\star$]
Quantitatively compare (BLEU, diversity) generation via greedy, beam search
($B=5$), and nucleus sampling ($p=0.9$) on a test corpus.
\end{exercise}

\begin{exercise}[$\star\star\star$]
Implement speculative decoding with a draft model (GPT-2 small) and a target
model (GPT-2 large). Measure the speedup obtained and verify that the output
distribution is preserved.
\end{exercise}
